%
% 14-integration.tex -- Integration der geometrischen Reihe
%
% (c) 2026 Prof Dr Andreas Müller
%
\subsection{Integration der geometrischen Reihe}
Durch Integration der geometrischen Reihe entlang einer Kurve
$\gamma$, kann eine auf dem Bild $\operatorname{im}\gamma = \gamma(I)$
definierte Funktion $f\colon\gamma(I)\to\mathbb{C}$ zu einer
Funktion erweitert werden, die automatisch holomorph und analytisch
ist.

\begin{satz}
Sie $\gamma\colon I\to\mathbb{C}$ eine differenzierbare Kurve und $g$
eine auf dem Bild $\operatorname{im}(\gamma)=\gamma(I)$ definierte 
stetige komplexe Funktion.
Dann ist die auf $\mathbb{C}\setminus\operatorname{im}\gamma$ definierte
Funktion
\begin{equation}
f(z)
=
\frac{1}{2\pi i}
\int_\gamma
\frac{g(x)}{x-z\mathstrut}\,dx
\end{equation}
holomorph.
\end{satz}

\begin{proof}
Da das Integrationsgebiet kompakt ist und $f$ stetig ist, kann unter
dem Integral differenziert werden.
Es ergibt sich die Ableitung
\begin{align*}
f'(z)
&=
\frac{1}{2\pi i}
\frac{d}{dz}
\int_\gamma
\frac{g(x)}{x-z\mathstrut}
\,dx
=
\frac{1}{2\pi i}
\int_\gamma
\frac{d}{dz}
\frac{g(x)}{x-z\mathstrut}
\,dx
=
\frac{1}{2\pi i}
\int_\gamma
\frac{g(x)}{(x-z\mathstrut)^2}
\,dx
\end{align*}
Insbesondere ist $f(z)$ komplex differenzierbar und damit $f$ holomorph.
\end{proof}

Andererseits kann man den Integranden mit der geometrischen Reihe
entwickeln und termweise integrieren.
Dazu verwenden wir
\eqref{buch:analytisch:eqn:geometrisch1}
und
\eqref{buch:analytisch:eqn:geometrischintegriert}
und erhalten
\begin{align}
f(z)
&=
\frac{1}{2\pi i}
\int_\gamma
\frac{g(x)}{x-a}
\cdot
\frac{1}{\displaystyle1-\frac{z-a\mathstrut}{x-a\mathstrut}}
\,dx
\notag
\\
&=
\frac{1}{2\pi i}
\int_\gamma
\frac{g(x)}{x-a}
\sum_{k=0}^\infty
\biggl(
\frac{z-a\mathstrut}{x-a\mathstrut}
\biggr)^k
\,dx
\notag
\\
&=
\sum_{k=0}^\infty
\biggl(
\frac{1}{2\pi i}
\int_\gamma
\frac{g(x)}{(x-a)^{k+1}}
\,dx
\biggr)
(z-a)^k.
\label{buch:analytisch:eqn:reihe}
\end{align}
Damit ist eine Potenzreihe im Punkt $a$ für die Funktion
$f(z)$ gefunden.
Sie ist konvergent, solange 
\[
\biggl|
\frac{z-a}{x-a}
\biggr|
=
\frac{|z-a|}{|x-a|}
<
1
\qquad\Rightarrow\qquad
|z-a| < |x-a|
\]
für alle $x$ entlang des Weges $\gamma$ gilt.
Dies wird erreicht, wenn 
\[
|z-a| < \inf_{\gamma} |x-a|,
\]
auf der rechten Seite steht der kleinste Abstand des Punktes $a$ vom
Weg $\gamma$.
Wir fassen dieses Resultat im folgenden Satz zusammen.

\begin{satz}
Sei $\gamma\colon I\to\mathbb{C}$ eine Kurve in $\mathbb{C}$ und
$f$ eine stetige Funktion auf dem Bild $\gamma(I)=\operatorname{im}\gamma$
der Kurve.
Dann ist die Funktion 
\[
f(z)
=
\int_{\gamma} \frac{g(x)}{x-z}\,dx
\]
analytisch und hat im Punkt $a\in\mathbb{C}\setminus \operatorname{im}\gamma$
die Potenzreihenentwicklung
\[
f(z)
=
\sum_{k=0}^\infty
\biggl(
\frac1{2\pi i}
\int_\gamma \frac{g(x)}{(x-a)^{k+1}}\,dx
\biggr)
(z-a)^k
\]
mit Konvergenzradius
\[
\varrho(a)
=
\inf\{ |\gamma(t)-a|\mid t\in I\}.
\]
\end{satz}
